(1) 用 \(c^2=gh\) 写 \(\partial_t(h,u)^\top+A\partial_x(h,u)^\top=0\)，其中 \(A=\begin{pmatrix}u&h\\g&u\end{pmatrix}\)，特征速度 \(\lambda_\pm=u\pm c\)。沿 \(dx/dt=u\pm c\) 的不变量为 \(J_\pm=u\pm2c\)。
(2) 连接左静水至右干床的是第一族稀疏扇：保持 \(J_+=u+2c=2c_0\)，令 \(\xi=\lambda_-=u-c\)。联立得
\[
u(\xi)=\frac23(c_0+\xi),\quad c(\xi)=\frac13(2c_0-\xi),\quad h(\xi)=\frac{(2c_0-\xi)^2}{9g}.
\]
完整解：\(\xi\le-c_0\) 时 \(h=h_0,u=0\)；\(-c_0<\xi<2c_0\) 时取上述表达式；\(\xi\ge2c_0\) 时 \(h=0\)，干床中的 \(u\) 无物理定义（不应强行给真空区赋速）。波尾速度 \(-c_0\)，湿润前沿 \(x_f=2c_0t\)，其水侧速度极限 \(u_f=2c_0\)，前沿随物质运动。两端深度连续，左端速度连续。
(3) 湿润的右半轴水质量（按单位宽度）为
\[
M_+(t)=\int_0^{2c_0t}h\,dx=\frac{t}{9g}\int_0^{2c_0}(2c_0-\xi)^2d\xi=\frac{8c_0^3}{27g}t.
\]
而 \(h(0,t)u(0,t)=(4c_0^2/9g)(2c_0/3)=8c_0^3/(27g)\)，故 \(\dot M_+=hu|_{x=0}\)。以连续方程消元将动量方程变成
\[
(hu)_t+\bigl(hu^2+\tfrac12gh^2\bigr)_x=0.
\]
在 \(x_f\) 有 \(h\to0\)、\(hu\to0\)、\(hu^2+gh^2/2\to0\)，故不存在有限的动量或质量脉冲逸出前沿。
(4) 波扇中 \(h=t^{-2}(2c_0t-x)^2/(9g)\)，且 \(u=2(c_0t+x)/(3t)\)。乘动量方程以 \(u\)，结合连续方程，得
\[
\partial_t\left(\tfrac12hu^2+\tfrac12gh^2\right)+
\partial_x\left[u\left(\tfrac12hu^2+gh^2\right)\right]=0.
\]
各项量纲为单位宽度的能量面密度除密度后乘速度；前沿 \(h\downarrow0\) 时能流趋零。设底摩擦或毛细力非零，则干床附近纵深趋零而相对作用增强，需另引入接触线/湿润前沿模型。
